\section{方法评价与阶段性结论}
\subsection{模型的适用性与局限}
本文框架把单点能力、调度、中继和分区组织在同一套物理与资源定义之下。其可解释性来自对象层次的区分：货箱是交付最小单位，架次是飞行任务单位，实体机与电池是分别占用的资源，通信分段是实际保障关系的记录单位。这样可以定位容量不足、能源等待或中继复制的具体来源。

模型的精确性范围也具有层次。单点集合划分在候选完整时可精确求解；固定时序区间模型可精确核算资源数；多点运输和静态中继候选的联合构造仍受启发式与候选集合限制。即使连续通信检验通过，也只证明该方案满足指定传播判据，并不证明中继点位或整体目标达到全局最优。

实际救援还受到天气、定位误差、临时禁飞、充电设备条件和动态需求影响。当前模型采用附件确定参数，30米DEM不能表征全部小尺度障碍，题设传播公式也不是完整无线环境仿真。因此，得到的方案应理解为标准场景下的计划依据，其推广需要额外的气象、通信和装备实测数据。

\subsection{主要结论}
单点任务中，B型仅在S008受到能量约束，C型在5个服务区的安全载荷低于额定值；完整货箱不可拆分又使部分批次先受到体积约束。最终单点运输需18架次。多点混装和资源事件调度将运输架次数降至22，并保证31个硬时限货箱全部按时送达。

通信约束使中继在位时间成为运输可行性的组成部分。当前方案以11个中继架次保障22个运输架次，连续区间认证和加密采样均未发现中断；相应等待使一般物资迟到指标明显增大，说明联合调度仍有压缩空间。固定该方案后，运输路线形成11个不可拆原子块。两组独立执行不产生库存缺口，三组执行则需要增加1架中继无人机；资源可实施性与组间工作量均衡之间存在明显取舍。
